% Section A.1, the solution to the hand calculation in Chapter 1, from the "Linear Regression by
% Hand" section of lectures/L01/exercises/README.md.

\section{Chapter 1: Linear Regression}
\label{sol:sec:c1}

\solution{c1:ex:1}{Training by hand}
Five training sets defined by $y = 3x + 2$:

\begin{narrowtable}{@{}cc@{}}
  $x$ & $y$ \\
  \midrule
  0 & 2 \\
  1 & 5 \\
  2 & 8 \\
  3 & 11 \\
  4 & 14 \\
\end{narrowtable}

The model starts out with both parameters at zero, and is trained for one epoch with a learning
rate of 10\%:
\[
  \begin{cases} k = 0 \\ m = 0 \end{cases}, \qquad \LR = 0.1
\]
The update rule is the one from \secref{c1:sec:linreg}, applied to every training set, including
the one at $x = 0$:
\[
  \delta = \yref - \yp, \quad \Delta e = \delta \cdot \LR, \quad m = m + \Delta e, \quad
  k = k + \Delta e \cdot x
\]
The training sets are used one at a time, in the order listed above.

\step{Training set 1}
From the first training set we get input $x = 0$ and reference value $\yref = 2$:
\[
  \begin{cases} x = 0 \\ \yref = 2 \end{cases}
\]
Both parameters are zero at the start, so the predicted output $\yp$ is zero as well:
\[
  \yp = k \cdot x + m = 0 \cdot 0 + 0 = 0
\]
The deviation $\delta$ is therefore two:
\[
  \delta = \yref - \yp = 2 - 0 = 2
\]
For a learning rate $\LR$ of 10\%, the adjustment amount $\Delta e$ is 0.2:
\[
  \Delta e = \delta \cdot \LR = 2 \cdot 0.1 = 0.2
\]
The m-value is increased directly by the adjustment amount $\Delta e$:
\[
  m = m + \Delta e = 0 + 0.2 = 0.2
\]
The k-value is increased by $\Delta e$ multiplied by $x$, which when $x = 0$ leaves it untouched:
\[
  k = k + \Delta e \cdot x = 0 + 0.2 \cdot 0 = 0
\]
After the first round of training:
\[
  \begin{cases} k = 0 \\ m = 0.2 \end{cases}
\]

\step{Training set 2}
\begin{align*}
  &\begin{cases} x = 1 \\ \yref = 5 \end{cases} \\
  \yp &= k \cdot x + m = 0 \cdot 1 + 0.2 = 0.2 \\
  \delta &= \yref - \yp = 5 - 0.2 = 4.8 \\
  \Delta e &= \delta \cdot \LR = 4.8 \cdot 0.1 = 0.48 \\
  m &= m + \Delta e = 0.2 + 0.48 = 0.68 \\
  k &= k + \Delta e \cdot x = 0 + 0.48 \cdot 1 = 0.48
\end{align*}
After the second round of training:
\[
  \begin{cases} k = 0.48 \\ m = 0.68 \end{cases}
\]

\step{Training set 3}
\begin{align*}
  &\begin{cases} x = 2 \\ \yref = 8 \end{cases} \\
  \yp &= k \cdot x + m = 0.48 \cdot 2 + 0.68 = 1.64 \\
  \delta &= \yref - \yp = 8 - 1.64 = 6.36 \\
  \Delta e &= \delta \cdot \LR = 6.36 \cdot 0.1 = 0.636 \\
  m &= m + \Delta e = 0.68 + 0.636 = 1.316 \\
  k &= k + \Delta e \cdot x = 0.48 + 0.636 \cdot 2 = 1.752
\end{align*}
After the third round of training:
\[
  \begin{cases} k = 1.752 \\ m = 1.316 \end{cases}
\]
The deviation has grown for three sets in a row, simply because the reference values grow faster
than the parameters have managed to so far.

\step{Training set 4}
\begin{align*}
  &\begin{cases} x = 3 \\ \yref = 11 \end{cases} \\
  \yp &= k \cdot x + m = 1.752 \cdot 3 + 1.316 = 6.572 \\
  \delta &= \yref - \yp = 11 - 6.572 = 4.428 \\
  \Delta e &= \delta \cdot \LR = 4.428 \cdot 0.1 = 0.4428 \\
  m &= m + \Delta e = 1.316 + 0.4428 = 1.7588 \\
  k &= k + \Delta e \cdot x = 1.752 + 0.4428 \cdot 3 = 3.0804
\end{align*}
After the fourth round of training:
\[
  \begin{cases} k = 3.0804 \\ m = 1.7588 \end{cases}
\]
The deviation has decreased for the first time, and the k-value has passed its target of 3.

\step{Training set 5}
\begin{align*}
  &\begin{cases} x = 4 \\ \yref = 14 \end{cases} \\
  \yp &= k \cdot x + m = 3.0804 \cdot 4 + 1.7588 = 14.0804 \\
  \delta &= \yref - \yp = 14 - 14.0804 = -0.0804
\end{align*}
The deviation is negative for the first time: the prediction overshoots the reference value, so
both parameters are adjusted downwards.
\begin{align*}
  \Delta e &= \delta \cdot \LR = -0.0804 \cdot 0.1 = -0.00804 \\
  m &= m + \Delta e = 1.7588 - 0.00804 = 1.75076 \\
  k &= k + \Delta e \cdot x = 3.0804 - 0.00804 \cdot 4 = 3.04824
\end{align*}
After the fifth and final round of training:
\[
  \begin{cases} k = 3.04824 \\ m = 1.75076 \end{cases}
\]
After a single epoch, the model predicts according to the following formula:
\[
  \yp = 3.04824 \cdot x + 1.75076
\]

\step{Prediction}
Predicted output $\yp$ and reference values $\yref$ for input $x$ in the range $[-5, 5]$, with the
predictions rounded to two decimal places:

\begin{narrowtable}{@{}rrr@{}}
  $x$ & $\yp$ & $\yref$ \\
  \midrule
  $-5$ & $-13.49$ & $-13$ \\
  $-4$ & $-10.44$ & $-10$ \\
  $-3$ & $-7.39$ & $-7$ \\
  $-2$ & $-4.35$ & $-4$ \\
  $-1$ & $-1.30$ & $-1$ \\
  0 & 1.75 & 2 \\
  1 & 4.80 & 5 \\
  2 & 7.85 & 8 \\
  3 & 10.90 & 11 \\
  4 & 13.94 & 14 \\
  5 & 16.99 & 17 \\
\end{narrowtable}

The weight is already within 2\% of its target after one epoch, while the bias is still 0.25
short, which is what tilts every prediction in the table. The bias converges more slowly here
because it's only ever adjusted by $\Delta e$, whereas the weight is adjusted by
$\Delta e \cdot x$, i.e.\ up to four times as much per training set.

The \code{x = 0} shortcut used by the code in Exercise~\ref{c1:ex:13} would have assigned $m = 2$
directly on the first training set, and the run would have ended somewhere else entirely. That's why
the exercise asks for the rule from \secref{c1:sec:linreg} here: it's the rule the hand calculations
in \secref{c1:sec:byhand} demonstrate, and it converges on both parameters from the training data
alone.
